% Options for packages loaded elsewhere
\PassOptionsToPackage{unicode}{hyperref}
\PassOptionsToPackage{hyphens}{url}
\documentclass[
  journal]{IEEEtran}
\usepackage{xcolor}
\usepackage{amsmath,amssymb}
\setcounter{secnumdepth}{-\maxdimen} % remove section numbering
\usepackage{iftex}
\ifPDFTeX
  \usepackage[T1]{fontenc}
  \usepackage[utf8]{inputenc}
  \usepackage{textcomp} % provide euro and other symbols
\else % if luatex or xetex
  \usepackage{unicode-math} % this also loads fontspec
  \defaultfontfeatures{Scale=MatchLowercase}
  \defaultfontfeatures[\rmfamily]{Ligatures=TeX,Scale=1}
\fi
\usepackage{lmodern}
\ifPDFTeX\else
  % xetex/luatex font selection
\fi
% Use upquote if available, for straight quotes in verbatim environments
\IfFileExists{upquote.sty}{\usepackage{upquote}}{}
\IfFileExists{microtype.sty}{% use microtype if available
  \usepackage[]{microtype}
  \UseMicrotypeSet[protrusion]{basicmath} % disable protrusion for tt fonts
}{}
\makeatletter
\@ifundefined{KOMAClassName}{% if non-KOMA class
  \IfFileExists{parskip.sty}{%
    \usepackage{parskip}
  }{% else
    \setlength{\parindent}{0pt}
    \setlength{\parskip}{6pt plus 2pt minus 1pt}}
}{% if KOMA class
  \KOMAoptions{parskip=half}}
\makeatother
\usepackage{longtable,booktabs,array}
\newcounter{none} % for unnumbered tables
\usepackage{calc} % for calculating minipage widths
% Correct order of tables after \paragraph or \subparagraph
\usepackage{etoolbox}
\makeatletter
\patchcmd\longtable{\par}{\if@noskipsec\mbox{}\fi\par}{}{}
\makeatother
% Allow footnotes in longtable head/foot
\IfFileExists{footnotehyper.sty}{\usepackage{footnotehyper}}{\usepackage{footnote}}
\makesavenoteenv{longtable}
\setlength{\emergencystretch}{3em} % prevent overfull lines
\providecommand{\tightlist}{%
  \setlength{\itemsep}{0pt}\setlength{\parskip}{0pt}}
\usepackage{bookmark}
\IfFileExists{xurl.sty}{\usepackage{xurl}}{} % add URL line breaks if available
\urlstyle{same}
% fallback for those not using the hyperref driver hyperxmp:
\makeatletter
\@ifundefined{xmpquote}{\newcommand{\xmpquote}[1]{#1}}{}
\makeatother
\hypersetup{
  hidelinks,
  pdfcreator={LaTeX via pandoc}}

\author{}
\date{}

\begin{document}

\section{Autonomous Navigation and Localization for Unmanned Aerial
Vehicles in GPS-Denied Environments: A Systematic Literature
Review}\label{autonomous-navigation-and-localization-for-unmanned-aerial-vehicles-in-gps-denied-environments-a-systematic-literature-review}

\textbf{Version 1 manuscript (V1; N = 171)}\\
Abhishek Raj

\subsection{Abstract}\label{abstract}

GPS/GNSS denial removes a primary source of absolute position and
exposes the limits of autonomous aerial navigation. This systematic
literature review synthesized methods for localization and navigation of
unmanned aerial vehicles (UAVs) in GPS-denied environments. The
Version-1 full-text corpus comprised 171 papers (26.9\% coverage of the
636-record included set); 465 included records were not retrieved.
Eligibility required a UAV-relevant platform, GPS/GNSS denial as the
primary contribution, quantitative experimental or simulated results,
peer review, publication from 2010-01-01 through 2026-06-30, and English
text. A CSV-driven extraction and 0--10 appraisal rubric were audited
against the Version-1 PDF set. Quality tiers were Q-high 38/171
(22.2\%), Q-medium 84/171 (49.1\%), and Q-low 49/171 (28.7\%); citation
tiers were Core 35/171 (20.5\%), Important 87/171 (50.9\%), and
Peripheral 49/171 (28.7\%). Experiments were Real\_World 22/171
(12.9\%), Simulation 78/171 (45.6\%), and Both 71/171 (41.5\%). The
synthesis indicates a progression from visual odometry and SLAM toward
visual--inertial, LiDAR--inertial, radio, and multi-sensor approaches,
while real-world and adversarial validation remain uneven. Findings are
limited by partial full-text coverage and machine-assisted extraction.

\textbf{Keywords:} UAV; GPS-denied; GNSS-denied; visual-inertial
odometry; LiDAR SLAM; sensor fusion; systematic literature review;
electronic warfare; localization.

\subsection{I. Introduction}\label{i-introduction}

\subsubsection{A. Operational
motivation}\label{a-operational-motivation}

UAV autonomy in indoor, urban, underground, forest, maritime, and
adversarial/electronic-warfare (EW) settings cannot assume continuous
GNSS. The review therefore treats denial, degradation, or jamming as a
core operating condition rather than a peripheral robustness test.

\subsubsection{B. Historical evolution}\label{b-historical-evolution}

The V1 corpus spans 2010--2026. The evidence narrative moves from indoor
visual odometry and early SLAM, through visual--inertial odometry (VIO),
to LiDAR--inertial fusion, radio and cooperative localization, and
recent learning- and geo-localization approaches. This is a qualitative
evolution statement; no unreported performance values are inferred.

\subsubsection{C. Gap and contributions}\label{c-gap-and-contributions}

Prior reviews often combine general navigation with GPS-denied use cases
or do not expose full-text coverage. This review contributes: (1) an
explicit primary-contribution eligibility test; (2) a reproducible CSV
schema and 0--10 quality rubric; (3) a versioned V1 corpus with audit
traceability; and (4) separate accounting for real, simulated, and mixed
validation.

\subsubsection{D. Organization}\label{d-organization}

Section II summarizes related work. Section III specifies protocol and
methods. Section IV reports V1 results. Sections V--VII discuss
implications, limitations, and conclusions. Appendices and the
supplementary package provide reproducibility records.

\subsection{II. Related Work}\label{ii-related-work}

The literature contains several recurring method families: visual
odometry and VIO; visual, LiDAR, and LiDAR--inertial SLAM; inertial and
filter-based navigation; radio/UWB and cooperative localization; map- or
geo-localization; and learning-assisted estimators. These families are
not mutually exclusive in the extracted records. The present review
differs by requiring GPS/GNSS denial to be the primary contribution and
by preserving UNKNOWN or NOT\_REPORTED values rather than imputing
metrics.

\subsection{III. Review Protocol and
Methods}\label{iii-review-protocol-and-methods}

\subsubsection{A. Research questions and
PICOC}\label{a-research-questions-and-picoc}

RQ1 asks which methods were proposed during 2010--2026. RQ2 asks which
sensor modalities and fusion strategies were used and how prevalence
changed. RQ3 asks where methods were validated and the real/simulation
ratio. RQ4 asks which open challenges remain, especially in
adversarial/EW settings.

The PICOC population was fixed-wing, rotary-wing, and hybrid UAVs; the
intervention was localization/navigation without GPS/GNSS; comparisons
were method families and time periods; outcomes included accuracy,
drift, computation, validation status, and sensor requirements; and
context covered indoor, urban, underground, forest, adversarial/EW, and
maritime environments.

\subsubsection{B. Information sources and
search}\label{b-information-sources-and-search}

IEEE Xplore and Scopus were the named databases. The search window was
2010-01-01 through 2026-06-30. Search strings, source logs, and
retrieval status are preserved in the S1--S10 supplementary package. No
PDFs were downloaded for this manuscript-generation task.

\subsubsection{C. Eligibility and
screening}\label{c-eligibility-and-screening}

All I1--I6 had to hold: UAV platform or explicitly transferable UAV
result; GPS/GNSS-denied navigation/localization as the primary
contribution; quantitative localization/navigation results;
peer-reviewed journal or conference publication; in-window publication;
and English language. Any E1--E8 condition excluded a record:
theoretical-only; incidental denial; GPS-augmented; out-of-scope
subject; non-UAV without transferability; non-peer-reviewed or
insufficiently retrievable; outside the window; or non-English. The I2
primary-contribution rule was applied strictly.

\subsubsection{D. Deduplication and corpus
versioning}\label{d-deduplication-and-corpus-versioning}

Deduplication used DOI-exact matching and title fuzzy matching at the
protocol threshold. The manuscript is V1 and uses only the 171-row
extracted corpus. The planned target corpus is not substituted into any
V1 result. Counts in this document were reconciled to
\texttt{02\_data\_processed/extracted\_master\_v2.csv} and
\texttt{06\_analysis/audit/V1\_audit\_20260917T124255Z.json}.

\subsubsection{E. Extraction and quality
appraisal}\label{e-extraction-and-quality-appraisal}

The extraction schema includes bibliographic fields, platform, sensors,
method category, environment, experiment type, metrics, application,
multi-agent status, notes, quality dimensions, quality tier, and
citation tier. QA uses rigor (0--4), reporting (0--3), baseline (0--2),
and reproducibility (0--1), summed to 0--10. Q-high is 8--10, Q-medium
5--7, and Q-low 0--4. Citation tiers are Core, Important, and
Peripheral. Multi-label sensor and environment fields are not mutually
exclusive.

\subsubsection{F. Audit and
reproducibility}\label{f-audit-and-reproducibility}

The audit reported 170 PASS and 1 PASS\_EXCEPTION. The exception was
REC\_1137: equation-dense text produced a 52.5\% short-token ratio;
manual review found no genuine extraction failure. The reproducibility
package identifies every source path, schema, quality row, audit status,
and seed-42 spot-check row.

\subsection{IV. Results}\label{iv-results}

\subsubsection{A. PRISMA flow and corpus
coverage}\label{a-prisma-flow-and-corpus-coverage}

The review flow yielded 636 stage-1 included records, of which 465 were
not retrieved for V1 and 171 full texts were extracted. Coverage was
171/636 = 26.9\%.

\subsubsection{B. Quality and citation
tiers}\label{b-quality-and-citation-tiers}

Q-high comprised 38/171 (22.2\%), Q-medium 84/171 (49.1\%), and Q-low
49/171 (28.7\%). Citation tiers were Core 35/171 (20.5\%), Important
87/171 (50.9\%), and Peripheral 49/171 (28.7\%). These are categorical
counts from the authoritative CSV, not estimates beyond V1.

\subsubsection{C. Experiment type and
validation}\label{c-experiment-type-and-validation}

Real\_World validation occurred in 22/171 (12.9\%), Simulation in 78/171
(45.6\%), and Both in 71/171 (41.5\%). Mixed validation was common
enough to warrant a separate category; it must not be collapsed into
either real or simulation.

\subsubsection{D. Methods, sensors, environments, and
applications}\label{d-methods-sensors-environments-and-applications}

Method, sensor, environment, platform, and application distributions are
provided as CSV-derived tables. Because sensor and environment fields
are multi-label, their category totals can exceed 171 and should not be
interpreted as mutually exclusive study counts. Metrics reported in the
source include ATE/RMSE, drift, success rate, and computational latency;
values marked NOT\_REPORTED remain missing.

\subsubsection{E. Audit result}\label{e-audit-result}

All 171 extracted records had an audit disposition: 170 PASS and 1
PASS\_EXCEPTION. The exception was manually reviewed as an extraction
artefact rather than a genuine failure. A separate 20-paper human
validation remains a required submission task and is not represented as
completed here.

\subsection{V. Discussion}\label{v-discussion}

\subsubsection{A. RQ1: method landscape}\label{a-rq1-method-landscape}

The V1 record set supports a heterogeneous method landscape rather than
a single dominant estimator. Visual--inertial and SLAM families coexist
with LiDAR--inertial fusion, filter-based methods, radio/UWB,
cooperative methods, and map-based approaches. Method labels should
therefore be read as primary categories assigned by the extraction
schema.

\subsubsection{B. RQ2: sensors and
fusion}\label{b-rq2-sensors-and-fusion}

The corpus emphasizes combinations of IMU with cameras and/or LiDAR,
while radio and other modalities appear in narrower roles. Fusion is
attractive because it can trade visual observability, geometric
structure, and inertial continuity, but sensor availability and
computation constrain deployment.

\subsubsection{C. RQ3: environments and
sim-to-real}\label{c-rq3-environments-and-sim-to-real}

Indoor and urban scenarios recur, with underground, forest, maritime,
and adversarial/EW settings represented unevenly. The 78 simulation-only
and 71 mixed records show substantial dependence on simulated or hybrid
evidence; the 22 real-world records provide the direct operational
subset.

\subsubsection{D. RQ4: open challenges and
EW}\label{d-rq4-open-challenges-and-ew}

Open challenges include resilience to jamming and spoofing, degraded
visual texture, dynamic scenes, map uncertainty, calibration drift,
compute and power limits, and reproducible adversarial benchmarks. EW
claims should distinguish denial, degradation, jamming, spoofing, and
general GNSS absence.

\subsubsection{E. Implications}\label{e-implications}

A credible deployment claim should state sensors, environment, baseline,
metrics, and failure modes together. The quality rubric and audit fields
make those reporting gaps visible, but they do not establish causal
superiority between method families.

\subsection{VI. Limitations}\label{vi-limitations}

First, V1 covers 171/636 included records (26.9\%); 465 records were not
retrieved. Second, extraction was machine-assisted and
keyword/regex-oriented, so categorical labels require human validation.
Third, one audit row required a documented PASS\_EXCEPTION. Fourth,
searches were limited to IEEE Xplore and Scopus and the English-language
window. Fifth, multi-label categories inflate marginal totals and do not
represent independent studies. Sixth, NOT\_REPORTED metrics were not
imputed, and no ATE or accuracy value is invented.

\subsection{VII. Conclusion}\label{vii-conclusion}

The V1 evidence base shows a broad, evolving field of GPS-denied UAV
navigation, with visual--inertial, SLAM, LiDAR--inertial, cooperative,
radio, and learning-assisted families. Validation remains split between
simulation and mixed studies, with fewer real-world records. The V1
findings are suitable as a traceable working manuscript, but coverage
and human validation gates must be addressed before submission.

\subsection{Data Availability}\label{data-availability}

The analysis artifacts use relative repository paths. The authoritative
extracted data are
\texttt{02\_data\_processed/extracted\_master\_v2.csv}; audit evidence
is \texttt{06\_analysis/audit/V1\_audit\_20260917T124255Z.json} and
\texttt{06\_analysis/audit/MANUAL\_REVIEW.md}.

\subsection{Conflicts and Funding}\label{conflicts-and-funding}

No funding or competing-interest statement was provided in the governing
files; these fields require author completion before submission.

\subsection{References}\label{references}

{[}1{]} Page MJ et al., ``The PRISMA 2020 statement,'' \emph{BMJ}, 2021,
Art. no. n71.\\
{[}2{]} Rethlefsen ML et al., ``PRISMA-S,'' \emph{Syst. Rev.}, 2021,
Art. no. 39.\\
{[}3{]} Source papers are listed in the repository CSV and reference
package; DOI-level bibliography formatting is a submission gate.

\subsection{Appendix A. PRISMA 2020 checklist (27
items)}\label{appendix-a-prisma-2020-checklist-27-items}

{\def\LTcaptype{none} % do not increment counter
\begin{tabular}{rlll}
\toprule
Item & Requirement & Location/evidence & Status \\
\midrule

\bottomrule

1 & Title identifies systematic review & Title & PASS \\
2 & Structured abstract & Abstract & PASS \\
3 & Rationale & I--II & PASS \\
4 & Objectives/RQs & III-A & PASS \\
5 & Eligibility criteria & III-C & PASS \\
6 & Information sources & III-B & PASS \\
7 & Search strategy & Supplement S2 & PASS \\
8 & Selection process & III-C & PASS \\
9 & Data collection process & III-E & PASS \\
10 & Data items & III-E & PASS \\
11 & Study risk-of-bias/appraisal & III-E & PASS \\
12 & Effect measures & III-E/IV-D & PASS \\
13 & Synthesis methods & IV and tables & PASS \\
14 & Reporting-bias assessment & Supplement S8; author completion &
PENDING \\
15 & Certainty assessment & Supplement S8; author completion &
PENDING \\
16 & Study selection results & IV-A & PASS \\
17 & Study characteristics & Tables II--VII & PASS \\
18 & Risk-of-bias results & IV-B & PASS \\
19 & Results of individual studies & Quality table & PASS \\
20 & Results of syntheses & IV & PASS \\
21 & Reporting biases & Supplement S8 & PENDING \\
22 & Certainty of evidence & Supplement S8 & PENDING \\
23 & Discussion interpretation & V & PASS \\
24 & Limitations & VI & PASS \\
25 & Conclusions & VII & PASS \\
26 & Registration/protocol & Data availability and protocol files &
PASS \\
27 & Support/funding/conflicts & Conflicts and Funding & PENDING \\
\end{tabular}
}

\subsection{Appendix B. Search strings}\label{appendix-b-search-strings}

Use the preserved IEEE Xplore and Scopus query exports in
\texttt{08\_docs/SUPPLEMENTARY\_V1.md} S2. Any final database-specific
syntax, date, and last-search timestamp must be completed by the
authors.

\subsection{Appendix C. Quality-control
audit}\label{appendix-c-quality-control-audit}

{\def\LTcaptype{none} % do not increment counter
\begin{tabular}{ll}
\toprule
Check & Status \\
\midrule

\bottomrule

V1 counts reconciled to CSV/audit & PASS \\
Coverage 171/636 = 26.9\% & PASS \\
Legacy values absent from new artifacts & PASS \\
Multi-label captions identify non-mutual exclusivity & PASS \\
ATE/accuracy not invented & PASS \\
Audit disposition 170 PASS + 1 PASS\_EXCEPTION & PASS \\
Manual review exception documented & PASS \\
20-paper human validation & PENDING \\
Figures 1--9 listed & PASS \\
PRISMA 27-item checklist present & PASS \\
PRISMA-S fields present & PASS \\
V1/V2 separation & PASS \\
\end{tabular}
}

\subsection{Appendix D. Version log}\label{appendix-d-version-log}

V1 uses the 171 full-text records available in the current working
corpus. The target expansion is documented separately and is not used in
V1 results.

\end{document}

